%\chapter{EPILOGO}
\begin{epilogo}

La crisis actual se explica por la incapacidad de aumentar el recaudo y la incapacidad de retornar a la senda de gasto vigente antes de la pandemia. La problemática, como se ha analizado a lo largo del texto, no es coyuntural sino estructural. El desafío de corregir esta situación excede las capacidades de cualquier Gobierno en solitario y requiere una solución de Estado. Para ello, sería necesario superar importantes barreras políticas, legales y constitucionales. Superarlas requerirá consenso nacional, pedagogía, capital político y conciencia sobre los sacrificios de corto plazo, que serán significativos, pero necesarios para generar beneficios a largo plazo en términos de mayor inversión privada, crecimiento económico y estabilidad financiera y macroeconómica.

El escenario en ausencia de consolidación fiscal es incierto. No es posible anticipar con precisión cuánto tiempo podría sostenerse el crecimiento de la deuda pública, pero es claro que, si la tendencia continúa, tarde o temprano el país enfrentará una corrección abrupta del gasto por la falta de liquidez o una inflación acelerada como forma de reducir el peso real de la deuda y financiar el funcionamiento del Estado. A medida que el Gobierno recurre a fuentes de financiamiento cada vez más costosas, aumenta el riesgo de que factores externos —como un cambio en las condiciones financieras internacionales o una pérdida de confianza de los inversionistas— frenen de forma repentina los flujos de inversión hacia los títulos de deuda pública. En ese caso, la escasez de recursos podría paralizar al Estado y conducir a una situación de impagos. Las consecuencias de un episodio de esta naturaleza serían inciertas y dependerían del contexto político en el que se materialice, pero implicarán el incumplimiento de compromisos con distintos grados de sensibilidad social, política y financiera. %Alternativamente, si el Banco de la República pierde su independencia, se podría desencadenar una aceleración de la inflación, la cual funciona como un impuesto regresivo y sin control democrático.


Afortunadamente, el país aún está a tiempo de corregir su desequilibrio fiscal e iniciar un proceso de consolidación ordenado. Sin embargo, este plan debe implementarse con urgencia, antes de que la deuda pública alcance niveles más difíciles de manejar y se intensifiquen las presiones estructurales sobre el gasto. Estas presiones -que ya se perfilan con claridad- provendrán del pago de intereses de la deuda pública, de las transferencias a las regiones a través del Sistema General de Participaciones (\gls{sgp}), de los beneficios asociados a diversos regímenes pensionales, del costo de mantener el sistema de salud actual, y del aumento reciente del gasto de personal. Cuanto más se postergue el ajuste, mayor será su costo en términos sociales, económicos y políticos.

El ajuste fiscal necesario para retornar a niveles moderados de deuda pública probablemente superará los cuatro puntos del PIB. Su materialización requerirá, con alta probabilidad, una combinación de reformas tributarias sucesivas, una reforma al sistema de salud orientada a la racionalización de costos, un esfuerzo por revertir el aumento reciente del gasto en personal y ajustes al sistema pensional que, desafortunadamente, implicarán una reducción de los beneficios actualmente otorgados. Adicionalmente, será indispensable garantizar que la Ley de Competencias que desarrolla el Acto Legislativo que aumentó las transferencias al \gls{sgp} sea plenamente compatible con la sostenibilidad fiscal. Al mismo tiempo, será necesario diseñar mecanismos que aseguren que los ingresos adicionales derivados de futuras reformas tributarias no sean automáticamente transferidos a las entidades territoriales a través del \gls{sgp}, sino que se orienten prioritariamente al cierre del déficit fiscal. Paralelo a este esfuerzo de consolidación, el país debe fortalecer el marco institucional que regula la gestión de las finanzas públicas, con el fin de asegurar un ciclo presupuestal coherente con el principio de sostenibilidad fiscal.

Una consolidación fiscal efectiva requerirá, inevitablemente, no solo de la reducción y contención del gasto, sino también un aumento significativo y sostenido del recaudo tributario. En este frente, el Estado colombiano ha mostrado una limitada capacidad para incrementar sus ingresos de forma permanente. Esto obliga a revisar con decisión el conjunto de beneficios tributarios, especialmente en el IVA, donde se concentra el mayor potencial de recaudo. El desafío en este punto es enorme, pues no solo exige diseñar mecanismos eficaces de compensación para los hogares más vulnerables que se vean afectados por el aumento del IVA, sino también un compromiso gubernamental creíble, democrático y pedagógico que permita evitar una nueva crisis de legitimidad como la que desencadenó el estallido social de 2021.

Finalmente, aunque escapa del alcance de este libro, es fundamental pensar en reformas estructurales que impulsen el crecimiento económico de largo plazo. Un ajuste fiscal, por sí solo, no será suficiente para garantizar la sostenibilidad de las finanzas públicas si no está acompañado de políticas que dinamicen la inversión, la productividad y la generación de empleo. Para ello, resulta clave avanzar en una agenda que incluya el fortalecimiento de la libre competencia, una reforma laboral que flexibilice la contratación formal, reformas al sistema financiero que permitan una asignación más eficiente del riesgo y una política de innovación y adopción tecnológica orientada a resultados, pero sin recurrir a beneficios tributarios. En ausencia de una estrategia integral de crecimiento, cualquier esfuerzo de consolidación fiscal corre el riesgo de ser fiscalmente insuficiente, socialmente insostenible y políticamente inviable.

El desafío fiscal que enfrenta Colombia no admite soluciones fáciles ni promesas de atajos. Pero tampoco es insuperable. Corregir el rumbo exigirá decisión política, sentido de responsabilidad histórica, pedagogía y un compromiso colectivo con el futuro del país. Por encima de todo, requerirá la construcción y el fortalecimiento de una identidad de Estado, en la que los impuestos no se perciban como una herramienta del gobierno de turno, sino como el mecanismo indispensable para el funcionamiento de un sociedad moderna y civilizada.

\end{epilogo}
